\documentclass[../main]{subfiles}
\externaldocument[main-]{../main}
\begin{document}

\subsection{\acl{BFR}}\label{sec_bag_federation_reduction}

\begin{definition}\label{def_bkg}
The \ac{BKG} of a federation $F$ is the bag union $\biguplus_{G \in \operatorname{codomain}(F)} G$ of all its members' \acp{KG}.
A triple $t$ thus has multiplicity $|\{\, G \in \operatorname{codomain}(F) : t \in G \,\}|$ meaning the number of members whose \ac{KG} contains it.
\end{definition}

We show that any union-based federation strategy is equivalent to executing the original query against a \emph{\ac{BKG}} as defined in \Cref{def_bkg}, making the query execution of a \ac{CQ} follow bag semantics rather than bag-set semantics of a \ac{UCQ}.
One can think of this as a virtual operator that aggregates the federation members into a \ac{BKG} and rewrites the query to be evaluated exclusively over it.
This is particularly useful given that \ac{UCQ} containment under bag and bag-set semantics is undecidable, as discussed in \Cref{sec:related_works_qc}.
We establish this using \ac{ESA} as the canonical proof vehicle, through the following propositions.

We first handle the case where federation members have no overlapping query-relevant knowledge.

\begin{proposition}\label{prop_no_overlap}
A query performed over a federation where each triple pattern can only be answered by one federation member yields the same result when performed over a \ac{BKG}.
\end{proposition}

\begin{proof}
Since each triple pattern $\mathit{tp}_j$ can only be answered by one member $\mathit{fm}_i$, every $\fmu$ in \Cref{eq_esa} collapses to a single $\freq_{\mathit{fm}_i}^{\mathit{tp}_j}$, leaving only join operations:
\[
\fmj(\freq_{\mathit{fm}_i}^{\mathit{tp}_1}, \dots, \freq_{\mathit{fm}_{n_{\mathrm{fm}}}}^{\mathit{tp}_{n_{\mathrm{tp}}}})
\]
where $n_{\mathrm{fm}}$ denotes the number of query-relevant federation members and $n_{\mathrm{tp}}$ the number of triple patterns in the query.
Given that no \texttt{UNION} are present, no duplicates arise from the federation structure.

In the \ac{BKG}, since no query-relevant triple appears in more than one member's \ac{KG}, each relevant triple appears exactly once.
The join evaluation therefore produces the same result bag.
\end{proof}

We now consider the opposite case, where all federation members host an identical \ac{KG}.

\begin{proposition}\label{prop_duplicate_fed_mem}
For a federation where all members host an identical \ac{KG}, the query produces the same results when evaluated over the federation's \ac{BKG}.
\end{proposition}

\begin{proof}
When all $m$ members host an identical \ac{KG}, every $\freq_{\mathit{fm}_i}^{\mathit{tp}_j}$ produces the same result for a given $\mathit{tp}_j$.

By the SPARQL definition of \texttt{UNION}, the multiplicity of the result bag equals the sum of the multiplicities of each operand's solution bag, so this produces $m$ copies of each result of $\freq_{\mathit{fm}_i}^{\mathit{tp}_j}$.
By definition of the \texttt{JOIN}, the multiplicity will be at most $m^3$ times the number of times solution mappings can be merged.
We say \emph{at most} because, depending on the query, some solution mappings may not merge; for a given query, however, evaluation over a \ac{BKG} and over a federation exhibits the same behavior.
% mu(u) = m, because we have m identical partial solution mapping if they can all merge together
% the multiplicity of each solution mapping will be m since we are dealing with duplicate
% Thus we have m time m time m

Let's suppose $\Omega_1$ being one set of $mu$

we know that the multiplcitiy of $operatorname{multiplicity}(\mu \mid \Omega)$ is 1 because we do not do projections,
and we are in the context of a set semantic database thus $\freq_{\mathit{fm}_i}^{\mathit{tp}_j}$ has a multiplicity of 1.
\begin{multline*}
\operatorname{multiplicity}(\mu \mid \Omega) = \sum_{i}^{m} \operatorname{multiplicity}(\mu \mid \Omega) = \sum_{i}^m 1 = $m$
\end{multline*}

for each union branch we have a multiplicity of $m$.

Let;s join 2 union branch

we know that the multiplicity of each branch $\operatorname{multiplicity}(\mu_u \mid \Omega_i)$ has a multiplicity of $m$ given we are joining the first
two branch of the euqation
\begin{align*}
\operatorname{multiplicity}(\mu \mid \texttt{JOIN}(\Omega_1, \Omega_2)) = \\
\sum_{(\mu_1, \mu_2) \in m(\mu)} \operatorname{multiplicity}(\mu_1 \mid \Omega_1) \cdot \operatorname{multiplicity}(\mu_2 \mid \Omega_2) = \\
\sum_{(\mu_1, \mu_2) \in m(\mu)} m \cdot m
\end{align*}

For the multiplicity of join let's first analyse this equation.

since in a branch of union there can only be the same result then we know that either no solution can be merged thus
$m(\mu)$ is equal to 0 or

\begin{align*}
\max(|m(\mu)|) = \max(\{(\mu_1, \mu_2) | \operatorname{merge}(\mu_1 \in \Omega_1, \mu_2 \in \Omega_2) = \mu\}) \\
= m^3
\end{align*}

because every solution mapping from $\Omega_1$ can be mapped to a solution mapping from  $\Omega_2$.

thus the multiplicity of the join become,
\begin{align*}
\max(\operatorname{multiplicity}(\mu \mid \texttt{JOIN}(\Omega_1, \Omega_2))) = \\
\max(\sum_{(\mu_1, \mu_2) \in m(\mu)} \operatorname{multiplicity}(\mu_1 \mid \Omega_1) \cdot \operatorname{multiplicity}(\mu_2 \mid \Omega_2)) = \\
\max(\sum_{(\mu_1, \mu_2) \in m(\mu)} m \cdot m) = m^5
\end{align*}

for subsequent join



In a \ac{BKG} each triple will have a multiplicity of $m$, and given the definition of \texttt{JOIN} the multiplicity of each solution mapping will be $m^3$.
Both evaluations therefore produce the same result bag.


for \ac{BKG}


\end{proof}

The general case combines both observations: partial knowledge overlap is handled by the same UNION mechanism, with multiplicities determined by how many members hold each piece of data.

\begin{proposition}\label{prop_esa_bkg}
A query performed using \ac{ESA} over a federation yields the same results as a query over a \ac{BKG}.
\end{proposition}

\begin{proof}
Query-irrelevant triples do not affect results.
For query-relevant knowledge, \Cref{eq_esa} introduces a \texttt{UNION} over all $m$ members for each triple pattern, distributing every triple pattern to each federation member to prevent result loss from joins requiring data from multiple members.
When a triple pattern can be answered by $k$ out of $m$ members, the \texttt{UNION}s retrieves $k$ times the triples matching the pattern.
By \Cref{def_bkg}, that same triple appears with multiplicity $k$ in the \ac{BKG}.
Generalizing \Cref{prop_duplicate_fed_mem}, the join multiplicity for a solution mapping $\mu$ is at most a factor $k$ instead of $m$; the \ac{BKG} therefore produces the same result bag.
\end{proof}

\begin{corollary}\label{cor_bfr}
Any FedQPL query plan that distributes triple patterns to multiple federation members via \texttt{UNION} produces results equivalent to evaluating the original query over a \ac{BKG}.
\end{corollary}

\begin{proof}
A FedQPL query plan either routes each triple pattern to exactly one member or routes some patterns to multiple members via \texttt{UNION}.
In the first case, the plan takes the join-only form of \Cref{prop_no_overlap} and no duplicates arise from federation structure.
In the second case, for each triple pattern $\mathit{tp}_i$ distributed to $k_i$ members, the \texttt{UNION} produces $k_i$ copies of each result.
By \Cref{def_bkg}, that triple appears with multiplicity $k_i$ in the \ac{BKG}.
The multiplicities therefore coincide in both evaluations.
\end{proof}

In \ac{ASSF}, the source selection algorithm determines the form of the FedQPL query plan.
When the plan contains \texttt{UNION} over multiple members for some triple pattern, \Cref{cor_bfr} applies and the evaluation follows bag semantics.
When the plan uses only joins, \Cref{prop_no_overlap} applies and bag-set semantics hold.
The \emph{\acl{BFR}} thus covers the full spectrum of union-based federation strategies.

\end{document}
